\section{Introducción}

El crecimiento de los asistentes de voz y los sistemas domóticos ha
popularizado el uso de modelos de lenguaje para interpretar comandos en
lenguaje natural (``prendé la luz del living'', ``bajá el aire a 20
grados''). La mayoría de las soluciones comerciales dependen de modelos de
gran escala alojados en la nube, lo que introduce tres problemas recurrentes
en el contexto de un hogar: latencia de red, costo recurrente por uso de
API, y dependencia de conectividad — un comando tan simple como apagar una
luz no debería fallar por una caída de internet.

Una alternativa creciente es el uso de Small Language Models (SLMs) —
modelos de lenguaje de entre unos cientos de millones y pocos miles de
millones de parámetros — ejecutados localmente en el propio hub domótico
(una Raspberry Pi, un mini-PC o un microcontrolador de gama alta). Trabajos
recientes muestran que estos modelos pueden alcanzar, en tareas acotadas, un
desempeño competitivo frente a modelos mucho más grandes \cite{gupta2025slms}.
Sin embargo, la mayor parte de la literatura sobre SLMs y asistentes
domóticos evalúa el inglés, y no está claro cómo se comportan estos modelos
frente a estructuras propias del español rioplatense (voseo, diminutivos,
expresiones relativas como ``bajale un toque'') ni cuál es el costo real en
latencia de ejecutarlos en hardware modesto.

Este trabajo aborda esa brecha mediante un estudio comparativo controlado.
Las contribuciones de este paper son: (1) un dataset de 32 comandos de
domótica en español rioplatense, organizados en seis categorías
lingüísticas, con una estructura de referencia (ground truth) de cinco
campos; (2) una evaluación empírica de cuatro SLMs open-source —
SmolLM2-360M-Instruct, Qwen2.5-0.5B-Instruct, Qwen2.5-1.5B-Instruct y
SmolLM2-1.7B-Instruct — ejecutados en CPU, midiendo coincidencia exacta,
exactitud por campo, validez del formato de salida y latencia; y (3) un
análisis cualitativo de los patrones de error que identifica una debilidad
sistemática, independiente del tamaño del modelo, en la interpretación de
comandos relativos o cualitativos sin valor numérico explícito.
